\chapter*{Symbol- und Abkürzungsverzeichnis}
\addcontentsline{toc}{chapter}{Symbol- und Abkürzungsverzeichnis}

\section*{Lateinische Formelzeichen}

\begin{longtable}{p{3cm}p{8cm}p{3cm}}
\textbf{Symbol} & \textbf{Bedeutung} & \textbf{Einheit} \\
\addlinespace[0.5em]
$A$ & Wärmeübertragungsfläche & $\mathrm{m^2}$ \\
$AFUDC$ & Zinsen während der Bauzeit & $\mathrm{EUR}$ \\
$C_\mathrm{BM}$ & Bare-Module-Kosten & $\mathrm{EUR}$ \\
$c_\mathrm{el}$ & Strompreis & $\mathrm{EUR\,kWh^{-1}}$ \\
$c_{\mathrm{N_2}}$ & spezifische Stickstoffbereitstellungskosten & $\mathrm{EUR\,kg_{N_2}^{-1}}$ bzw. $\mathrm{EUR\,t_{N_2}^{-1}}$ \\
$c_p$ & spezifische Wärmekapazität & $\mathrm{kJ\,kg^{-1}\,K^{-1}}$ \\
$CC_L$ & nivellierte kapitalgebundene Kosten & $\mathrm{EUR\,a^{-1}}$ \\
$CEPCI$ & Chemical Engineering Plant Cost Index & $-$ \\
$CELF$ & Constant Escalation Levelization Factor & $-$ \\
$CRF$ & Kapitalrückgewinnungsfaktor & $-$ \\
$DC$ & direkte Kosten & $\mathrm{EUR}$ \\
$\bar{e}^{CH}_i$ & molare chemische Standardexergie der Komponente $i$ & $\mathrm{kJ\,mol^{-1}}$ \\
$\dot{E}$ & Exergiestrom & $\mathrm{W}$ \\
$EC_L$ & nivellierte Energiekosten & $\mathrm{EUR\,a^{-1}}$ \\
$FCI$ & Fixed Capital Investment & $\mathrm{EUR}$ \\
$f$ & Berechnungsfaktor & $-$ \\
$g$ & Erdbeschleunigung & $\mathrm{m\,s^{-2}}$ \\
$H$ & Enthalpie & $\mathrm{kJ\,mol^{-1}}$ \\
$h$ & spezifische Enthalpie & $\mathrm{kJ\,kg^{-1}}$ \\
$IC$ & indirekte Kosten & $\mathrm{EUR}$ \\
$i_\mathrm{eff}$ & effektiver Zinssatz & $-$ \\
$k_j$ & Kostenfaktor der Kostenart $j$ & $-$ \\
$LRD$ & Lizenz-, Forschungs- und Entwicklungskosten & $\mathrm{EUR}$ \\
$m_{N_2,a}$ & jährlich erzeugte Stickstoffmasse & $\mathrm{t\,a^{-1}}$ \\
$\dot{m}$ & Massenstrom & $\mathrm{kg\,s^{-1}}$ \\
$n$ & ökonomische Lebensdauer & $\mathrm{a}$ \\
$\dot{n}$ & Stoffmengenstrom & $\mathrm{mol\,s^{-1}}$ bzw. $\mathrm{kmol\,s^{-1}}$ \\
$N_\mathrm{theo}$ & Anzahl theoretischer Trennstufen & $-$ \\
$OM_L$ & nivellierte Betriebs- und Wartungskosten & $\mathrm{EUR\,a^{-1}}$ \\
$p$ & Druck & $\mathrm{bar}$ \\
$PEC$ & Purchased Equipment Costs & $\mathrm{EUR}$ \\
$\dot{Q}$ & Wärmestrom & $\mathrm{W}$ \\
$R$ & universelle Gaskonstante & $\mathrm{kJ\,mol^{-1}\,K^{-1}}$ \\
$r_{n,j}$ & nominale Eskalationsrate der Kostenart $j$ & $-$ \\
$s$ & spezifische Entropie & $\mathrm{kJ\,kg^{-1}\,K^{-1}}$ \\
$SUC$ & Anfahrkosten & $\mathrm{EUR}$ \\
$T$ & Temperatur & $\mathrm{K}$ \\
$TCI$ & Total Capital Investment & $\mathrm{EUR}$ \\
$t_\mathrm{op}$ & jährliche Betriebsstunden & $\mathrm{h\,a^{-1}}$ \\
$TRR_L$ & nivellierter Total Revenue Requirement & $\mathrm{EUR\,a^{-1}}$ \\
$U$ & innere Energie & $\mathrm{J}$ \\
$V$ & Volumen & $\mathrm{m^3}$ \\
$v$ & Geschwindigkeit & $\mathrm{m\,s^{-1}}$ \\
$\dot{V}$ & Volumenstrom & $\mathrm{m^3\,s^{-1}}$ bzw. $\mathrm{Nm^3\,h^{-1}}$ \\
$W$ & Arbeit & $\mathrm{J}$ \\
$\dot{W}$ & Arbeits- bzw. Leistungsstrom & $\mathrm{W}$ \\
$WC$ & Umlaufkapital & $\mathrm{EUR}$ \\
$x$ & Dampfqualität bzw. Stoffmengenanteil in der Flüssigphase & $-$ \\
$x_i$ & Molfraktion der Komponente $i$ & $-$ \\
$y$ & Stoffmengenanteil in der Gasphase & $-$ \\
$y_{D,b}$ & Exergievernichtungsverhältnis des Funktionsblocks $b$, bezogen auf die Brennstoffexergie der Gesamtanlage & $-$ \\
$y^*_{D,b}$ & relative Exergievernichtung des Funktionsblocks $b$, bezogen auf die gesamte Exergievernichtung der Anlage & $-$ \\
$z$ & Höhenkoordinate & $\mathrm{m}$ \\
\end{longtable}

\section*{Griechische Formelzeichen}

\begin{longtable}{p{3cm}p{8cm}p{3cm}}
\textbf{Symbol} & \textbf{Bedeutung} & \textbf{Einheit} \\
\addlinespace[0.5em]
$\Delta e$ & spezifische Exergiedifferenz & $\mathrm{kJ\,kg^{-1}}$ \\
$\Delta p$ & Druckverlust & $\mathrm{bar}$ bzw. $\mathrm{kPa}$ \\
$\Delta_\mathrm{R} H^\circ$ & Standard-Reaktionsenthalpie & $\mathrm{kJ\,mol^{-1}}$ \\
$\Delta T$ & Temperaturdifferenz & $\mathrm{K}$ \\
$\Delta T_{\min}$ & minimaler Temperaturabstand & $\mathrm{K}$ \\
$\varepsilon$ & exergetischer Wirkungsgrad & $-$ \\
$\eta$ & Wirkungsgrad & $-$ \\
$\eta_\mathrm{is}$ & isentroper Wirkungsgrad & $-$ \\
$\eta_\mathrm{mech}$ & mechanischer Wirkungsgrad & $-$ \\
$\eta_{\mathrm{N_2}}$ & Stickstoffausbeute & $-$ \\
$\rho$ & Dichte & $\mathrm{kg\,m^{-3}}$ \\
$\varphi$ & relative Luftfeuchte & $-$ \\
\end{longtable}

\section*{Tiefgestellte Indizes}

\begin{longtable}{p{3cm}p{11cm}}
\textbf{Index} & \textbf{Bedeutung} \\
\addlinespace[0.5em]
$0$ & Referenzzustand \\
$a$ & jährlich \\
$\mathrm{aus}$ & austretend \\
$b$ & Funktionsblock \\
$\mathrm{BM}$ & Bare Module \\
$\mathrm{BM,tot}$ & gesamte Bare-Module-Kosten \\
$\mathrm{D}$ & Destruction bzw. Exergievernichtung \\
$el$ & elektrisch \\
$\mathrm{ein}$ & eintretend \\
$ex$ & exergetisch \\
$\mathrm{F}$ & Fuel bzw. Brennstoffexergie \\
$ges$ & gesamt \\
$\mathrm{HP}$ & Hochdruck \\
$i$ & Komponente \\
$\mathrm{in}$ & Eintritt \\
$\mathrm{is}$ & isentrop \\
$j$ & Stoffstrom bzw. Kostenart \\
$k$ & Komponente \\
$L$ & Levelized bzw. annualisiert; bei Exergieströmen Loss bzw. Verlust \\
$\mathrm{LP}$ & Niederdruck \\
$\mathrm{Loop}$ & Synthesekreislauf \\
$\max$ & maximal \\
$\mathrm{mech}$ & mechanisch \\
$\min$ & minimal \\
$\mathrm{mix}$ & Mischzustand \\
$\mathrm{N}$ & Normbedingung \\
$\mathrm{O_2}$ & Sauerstoff \\
$\mathrm{OM}$ & Betriebs- und Wartungskosten \\
$\mathrm{out}$ & Austritt \\
$\mathrm{P}$ & Product bzw. Produktexergie \\
$\mathrm{pot}$ & potenziell \\
$\mathrm{R}$ & Reaktion \\
$\mathrm{S}$ & Skalierung \\
$\mathrm{sys}$ & System \\
$\mathrm{th}$ & thermisch \\
$\mathrm{theo}$ & theoretisch \\
$\mathrm{tot}$ & total bzw. gesamt \\
\end{longtable}

\section*{Hochgestellte Kennzeichnungen}

\begin{longtable}{p{3cm}p{11cm}}
\textbf{Kennzeichnung} & \textbf{Bedeutung} \\
\addlinespace[0.5em]
$\mathrm{CH}$ & chemisch \\
$\mathrm{KN}$ & kinetisch \\
$\mathrm{M}$ & mechanisch \\
$\mathrm{PH}$ & physikalisch \\
$\mathrm{PT}$ & potenziell \\
$\mathrm{T}$ & thermisch \\
$\circ$ & Standardzustand \\
\end{longtable}

\section*{Abkürzungen}

\begin{longtable}{p{3cm}p{11cm}}
\textbf{Abkürzung} & \textbf{Bedeutung} \\
\addlinespace[0.5em]
AFUDC & Allowance for Funds Used During Construction (Zinsen während der Bauzeit) \\
APSA & kompakte Stickstoffanlage von Linde \\
ASU & Air Separation Unit (Luftzerlegungsanlage) \\
ATR & Autothermal Reforming (autotherme Reformierung) \\
BM & Bare Module \\
CAPEX & Capital Expenditures (Investitionsausgaben) \\
CBM & Bare-Module-Kosten \\
CC & Capital Costs (kapitalgebundene Kosten) \\
CCS & Carbon Capture and Storage (CO\textsubscript{2}-Abscheidung und -Speicherung) \\
CEPCI & Chemical Engineering Plant Cost Index \\
CELF & Constant Escalation Levelization Factor (Nivellierungsfaktor) \\
CRF & Capital Recovery Factor (Kapitalrückgewinnungsfaktor) \\
CRYOSS & kryogene Stickstoffanlage von Linde \\
DC & Direct Costs (direkte Kosten) \\
DSM & Demand-Side-Management (Lastmanagement) \\
DWC & Dividing Wall Column (Trennwandkolonne) \\
EC & Energy Costs (Energiekosten) \\
EUR & Euro \\
FCI & Fixed Capital Investment (Anlageinvestitionskosten) \\
HI-stages & Heat-Integrated Stages (wärmeintegrierte Stufen) \\
HP & High Pressure (Hochdruck) \\
IC & Indirect Costs (indirekte Kosten) \\
KOL & Rektifikationskolonne des Einkolonnenmodells \\
KOLHP & Hochdruckkolonne des Doppelkolonnenmodells \\
KOLLP & Niederdruckkolonne des Doppelkolonnenmodells \\
LAES & Liquid Air Energy Storage (Flüssigluft-Energiespeicherung) \\
LP & Low Pressure (Niederdruck) \\
LPC & Low Pressure Column (Niederdruckkolonne) \\
LRD & Licensing, Research and Development Costs (Lizenz-, Forschungs- und Entwicklungskosten) \\
LV & Luftverdichtung \\
MPC & Model Predictive Control (modellprädiktive Regelung) \\
MSR & Mess- und Regelungstechnik \\
OM & Operation and Maintenance Costs (Betriebs- und Wartungskosten) \\
OPEX & Operating Expenditures (Betriebsausgaben) \\
PEC & Purchased Equipment Costs (Anschaffungskosten der Apparate) \\
ppm & parts per million (Teile pro Million) \\
PRISM & kompakte Stickstoffanlage von Air Products \\
PSC & Prepurified Single Column (Luftzerlegungsanlage mit vorgeschalteter Reinigung) \\
PSA & Pressure Swing Adsorption (Druckwechseladsorption) \\
RadFrac & Aspen-Plus-Modell für Rektifikationskolonnen \\
SMR & Steam Methane Reforming (Dampfreformierung) \\
SPECO & Specific Exergy Costing \\
SUC & Start-up Costs (Anfahrkosten) \\
SZ & zusätzlicher interner Stoffstrom zur Abbildung gekoppelter Wärmeübertrager \\
TCI & Total Capital Investment (Gesamtinvestition) \\
TCN-BE & kompakte Stickstoffanlage von Linde \\
TNSC & Taiyo Nippon Sanso Corporation \\
TRR & Total Revenue Requirement (jährlich erforderliche Gesamtkosten) \\
TSA & Temperature Swing Adsorption \\
USD & US-Dollar \\
VSA & Vacuum Swing Adsorption (Vakuumwechseladsorption) \\
WC & Working Capital (Umlaufkapital) \\
\end{longtable}

